\section{Discussion}
\label{sec:discussion}

A random split of this trace rewards memorisation: at most \V{full.random.maxnew} of the \V{full.split.random.s42.ntest} test tasks of a split come from unseen users, and \V{rf.full.dup.random.min} duplicate a training task exactly. The interaction models gain most there, and held-out users remove their lead. The pattern is consistent with that gain being knowledge of particular users and duplicated requests; whether interactions themselves fail to transfer would need a regularised or interaction-constrained boosting model, which we did not test. For known users the features carry over to a later period only weakly (\V{full.timeknown.minauc} to \V{full.timeknown.maxauc}); for new users they carry a weak signal at best. A deployment could at most support operators reviewing known users' recurring pipelines and should not judge a new user's jobs. The audit chain used only additivity, so the additive EBM would support it as well; the KAN form adds smooth curves with an explicit penalty. Compared with the first version, \kaam{} no longer matches the black-box baselines on random splits, and the single below-chance grouped partition of that version is one draw from a wide spread of fold results.
